\begin{frame}
  \frametitle{Einführung in SQL}
  \begin{itemize}
    \item SQL: Structured Query Language
    \item Bedeutung von Datenbanken
    \item Motivation für den Einsatz von SQL
  \end{itemize}
\end{frame}

\section{Grundlagen von SQL}
\begin{frame}[fragile]
  \frametitle{Grundlagen von SQL}
  \begin{itemize}
    \item SQL-Befehle für Tabellenerstellung
  \end{itemize}
  
  \begin{lstlisting}[language=SQL]
  CREATE TABLE Students (
    StudentID INT PRIMARY KEY,
    FirstName VARCHAR(50),
    LastName VARCHAR(50),
    Age INT
  );
  \end{lstlisting}
\end{frame}

\begin{frame}[fragile]
  \frametitle{Einfache SELECT-Anfrage}
  
  \begin{lstlisting}[language=SQL]
  SELECT FirstName, LastName
  FROM Students
  WHERE Age > 18;
  \end{lstlisting}
  
  \begin{table}[h]
    \centering
    \begin{tabular}{|c|c|}
      \hline
      FirstName & LastName \\
      \hline
      John & Doe \\
      Jane & Smith \\
      \hline
    \end{tabular}
    \caption{Ergebnis der SELECT-Anfrage}
  \end{table}
\end{frame}


% Fortsetzung der Präsentation
\begin{frame}[fragile]
  \frametitle{Daten einfügen (INSERT)}
  
  \begin{lstlisting}[language=SQL]
  INSERT INTO Students (FirstName, LastName, Age)
  VALUES ('Alice', 'Johnson', 20);
  \end{lstlisting}
\end{frame}

\begin{frame}[fragile]
  \frametitle{Daten aktualisieren (UPDATE)}
  
  \begin{lstlisting}[language=SQL]
  UPDATE Students
  SET Age = 21
  WHERE LastName = 'Johnson';
  \end{lstlisting}
\end{frame}

\begin{frame}[fragile]
  \frametitle{Daten löschen (DELETE)}
  
  \begin{lstlisting}[language=SQL]
  DELETE FROM Students
  WHERE Age < 20;
  \end{lstlisting}
\end{frame}

\section{Fortgeschrittene SQL-Abfragen}
\begin{frame}[fragile]
  \frametitle{JOIN-Operationen}
  
  \begin{lstlisting}[language=SQL]
  SELECT Students.FirstName, Students.LastName, Courses.CourseName
  FROM Students
  INNER JOIN Enrollments ON Students.StudentID = Enrollments.StudentID
  INNER JOIN Courses ON Enrollments.CourseID = Courses.CourseID;
  \end{lstlisting}
\end{frame}

\section{Datensicherheit}
\begin{frame}[fragile]
  \frametitle{Datensicherheit und SQL Injection}
  \begin{itemize}
    \item Parameterisierte Abfragen verwenden
  \end{itemize}
  
  \begin{lstlisting}[language=SQL]
  $sql = "SELECT * FROM Users WHERE Username = ? AND Password = ?";
  $stmt = $conn->prepare($sql);
  $stmt->bind_param("ss", $username, $password);
  $stmt->execute();
  \end{lstlisting}
\end{frame}

\section{Zusammenfassung}
\begin{frame}
  \frametitle{Zusammenfassung}
  \begin{itemize}
    \item Wichtige SQL-Befehle
    \item SELECT, INSERT, UPDATE, DELETE
    \item Fortgeschrittene Abfragen mit JOIN
    \item Datensicherheit und SQL Injection-Prävention
  \end{itemize}
\end{frame}

\section{Fragen}
\begin{frame}
  \frametitle{Fragen und Diskussion}
  Vielen Dank für Ihre Aufmerksamkeit! \\
  Haben Sie Fragen oder möchten Sie etwas diskutieren?
\end{frame}